\documentclass[11pt]{article}
\usepackage{amsmath,amsthm,amssymb,booktabs,listings}
\lstset{basicstyle=\ttfamily\small, breaklines=true}
\newtheorem{theorem}{Theorem}
\newtheorem{proposition}[theorem]{Proposition}

\title{N=12 Exhaustive EML Search: Plan and Progress}
\author{A.R.\ Almaguer}
\date{2026-04-20}

\begin{document}
\maketitle

\begin{abstract}
We document the plan and current progress for an exhaustive search over all EML trees
with up to $N=12$ nodes. The search aims to: (a) certify that no function representable
in $\leq 11$ nodes is representable in fewer, (b) find novel 12-node functions not in
the current SuperBEST table, and (c) discover potential new theorems about EML depth limits.
\end{abstract}

\section{Search Space}

An EML tree of $n$ nodes is a full binary tree where:
\begin{itemize}
\item Each internal node computes one of the 16 EML-family operators
\item Each leaf is a variable ($x_1, \ldots, x_k$) or a rational constant
\item The tree has $n$ internal nodes and $n+1$ leaves
\end{itemize}

\textbf{Search space size at N nodes:}
\begin{align}
|\mathcal{T}_n| &= C_n \cdot 16^n \cdot (k+m)^{n+1}
\end{align}
where $C_n = \frac{1}{n+1}\binom{2n}{n}$ is the $n$-th Catalan number (tree shapes),
$k$ = number of input variables, $m$ = number of constant leaves.

\begin{table}[h]
\centering
\caption{Search space sizes.}
\begin{tabular}{@{}rccc@{}}
\toprule
$n$ & $C_n$ & $16^n$ & Approx.\ total (2 vars, 3 consts) \\
\midrule
1  & 1      & 16     & $\sim 400$ \\
2  & 2      & 256    & $\sim 32{,}000$ \\
3  & 5      & 4096   & $\sim 10^7$ \\
6  & 132    & $\sim 10^7$ & $\sim 10^{15}$ \\
9  & 4862   & $\sim 10^{11}$ & $\sim 10^{21}$ \\
12 & 208{,}012 & $\sim 10^{14}$ & $\sim 10^{28}$ \\
\bottomrule
\end{tabular}
\end{table}

Full exhaustive search at $N=12$ is computationally infeasible ($\sim 10^{28}$).
The strategy is \textbf{filtered search} using:

\section{Search Strategy}

\subsection{Phase 1: Symbolic normalization (already complete, N=11)}

Many trees are equivalent under:
\begin{itemize}
\item Commutativity: $\mathrm{eml}(a,b) \neq \mathrm{eml}(b,a)$ (EML is not commutative)
\item Constant folding: subtrees with no variable leaves reduce to a single constant
\item Identity elimination: $\mathrm{exl}(0,x) = \ln(x)$, $\mathrm{eml}(x,1) = \exp(x)$
\end{itemize}

After normalization, the effective search space shrinks by $\sim 100\times$.

\subsection{Phase 2: Function fingerprinting}

Each candidate tree is evaluated at a fixed set of test points
$\{(x_i, y_j)\} = \{0.5, 1.0, 1.5, 2.0, 2.5\}^2$.
Two trees with identical fingerprints are considered equivalent (probabilistic certificate).
This reduces the search to unique function classes.

\subsection{Phase 3: Beam search + MCTS (current)}

The Rust binary \texttt{monogate-core/src/bin/n12\_search.rs}:
\begin{lstlisting}
cargo run --release --bin n12_search -- \
  --max-nodes 12 \
  --operators eml,exl,eal,edl,deml \
  --target-functions "sin,cos,tan,sqrt,cbrt" \
  --beam-width 10000 \
  --output results/n12_candidates.json
\end{lstlisting}

\subsection{Phase 4: Near-miss analysis}

Functions that achieve relative error $< 10^{-4}$ on the test domain but are not exact
are flagged as \emph{near-misses}. These are candidates for:
\begin{itemize}
\item New depth-12 approximations (PROPOSITION tier if error bound is proved)
\item Near-miss obstruction theory (planned paper: \texttt{papers/near\_miss\_obstruction/})
\end{itemize}

\section{Current Results}

\begin{proposition}[Depth-6 Phase Transition — T25]
The first $\mathrm{EML}_1$ values with $\mathrm{Im} > 0$ appear at depth 6.
For depths 1--5, all EML trees with real inputs produce real outputs.
\end{proposition}

\begin{proposition}[N=11 Complete at N=9]
SuperBEST FINAL uses at most 3 nodes for all basic arithmetic operations.
The full N=9 space has been searched; no sub-3-node representation of multiplication
or addition exists (confirming T08, T10 optimality).
\end{proposition}

\section{Planned Outputs}

Upon completion:
\begin{enumerate}
\item \texttt{python/results/n12\_candidates.json} --- all novel 12-node functions found
\item New theorem candidates (if any new elementary function gets a shorter tree)
\item Near-miss catalog for obstruction paper
\item Updated SuperBEST table (if any entry improves)
\end{enumerate}

\section{Timeline}

\begin{itemize}
\item N $\leq$ 9: Complete (2026-04-20)
\item N = 10--11: Beam search in progress
\item N = 12: Planned (pending compute)
\item Full exhaustive N $\leq$ 6: Complete via Python \texttt{exhaustive\_search.py}
\end{itemize}

\end{document}
